\subsection{Regresión lineal}
\label{sec:regresion_lineal}
El modelo de regresión lineal es uno de los modelos más simples utilizados en el modelado de sistemas. Este modelo provee una estimación de una variable continua $y$ a partir de una variable independiente $x$, cuya relación se modela con una recta, como se muestra en la \refeq{linear_regression} \cite{Bishop2006}.

\begin{equation}
    \label{eq:linear_regression}
    \hat{y} = \theta_0 + \theta_1 x
\end{equation}

En el caso general de la regresión lineal, se busca modelar la relación entre una variable dependiente $y$ y una o más variables independientes $x_j$ mediante una combinación lineal de estas últimas, como se muestra en la \refeq{linear_regression_general}.

\begin{equation}
    \label{eq:linear_regression_multivariable}
    \hat{y} = \theta_0 + \sum_{j=1}^{m} \theta_j x_j
\end{equation}

Una elección común para la métrica utilizada en problemas de regresión lineal es el error cuadrático medio explicado en la sección \ref{sec:metricas_regresion}, definida en la \refeq{mse}. En este caso, la función de costo $\mathcal{L}$ a minimizar es directamente la métrica MSE, como se muestra en la \refeq{mse_regression}.

\begin{equation}
    \label{eq:mse_regression}
    \mathcal{L}(\theta, \mathcal{D}) = \frac{1}{n} \sum_{i=1}^{n} (y_i - \hat{y}_i)^2
\end{equation}

Para este problema particular, es posible deducir analíticamente los valores de $\mathbf{\theta}$ que minimizan la función de costo definida en la \refeq{mse_regression}.

Si se deriva la función de costo con respecto a los parámetros y se iguala a cero como se muestra en la \refeq{mse_derivative}.

\begin{equation}
    \label{eq:mse_derivative}
    \frac{\partial \mathcal{L}}{\partial \theta} = 0
\end{equation}

Se obtiene la fórmula conocida como OLS (por sus siglas en inglés \textit{Ordinary Least Squares}), que permite calcular los parámetros óptimos $\hat{\theta}$ de manera directa, como se muestra en la \refeq{ols}, donde $X$ es la matriz de características, $\hat{\theta}$ es el vector de parámetros, y $y$ es el vector de salidas reales, como se muestra en la \refeq{ols_matrices} \cite{Bishop2006}

\begin{equation}
    \label{eq:ols}
        \hat{\theta} = \left(X^T X\right)^{-1} X^T y
\end{equation}

\begin{equation}
    \label{eq:ols_matrices}
      X =
        \begin{bmatrix}
        1 & x_{1,1} & x_{2,1} & \dots & x_{n,1} \\
        1 & x_{1,2} & x_{2,2} & \dots & x_{n,2} \\
        \vdots & \vdots & \vdots & \dots & \vdots \\
        1 & x_{1,m} & x_{2,m} & \dots & x_{n,m}
        \end{bmatrix}
      \quad \quad
      \hat{\theta} =
        \begin{bmatrix}
          \hat{\theta}_0 \\ \hat{\theta}_1 \\ \hat{\theta}_2 \\ \dots \\ \hat{\theta}_n
        \end{bmatrix}
      \quad \quad
      y = \begin{bmatrix} y_1 \\ y_2 \\ \vdots \\ y_m \end{bmatrix}
    \end{equation}

Este problema es un caso de un problema de optimización con solución cerrada, ya que permite calcular los parámetros óptimos de manera directa sin necesidad de utilizar métodos iterativos.

A modo de ilustración, el resultado de aplicar este modelo a un conjunto de datos se muestra en la \reffig{linear_regression}, donde se observa la línea de regresión ajustada a los datos, así como los errores entre las predicciones del modelo y los datos reales.

\defaultfigurecustomsize{regression/linear_regression.png}{Regresión Lineal.}{linear_regression}{1}
